
\subsubsection{2.1 Constrained Decoding for Code Generation}

Grammar-based constrained decoding methods enforce hard syntactic constraints during generation to guarantee validity. SYNCHROMESH uses formal grammars to restrict the search space, ensuring 100\% syntactically valid outputs by only generating tokens that maintain grammar compliance. NeuroLogic Decoding and GeLM extend this approach to handle semantic constraints beyond pure syntax. While these methods achieve perfect syntactic validity, they introduce substantial computational overhead: grammar parsing at each generation step slows inference, and hard constraints can reduce generation quality by forcing the model into grammatically valid but semantically awkward constructions. The approach presented here differs---rather than enforcing constraints that guarantee validity, syntax is treated as a soft scoring signal ($\beta =0.3$ weight) that guides but does not dictate generation, achieving 76.22\% validity at <0.05ms AST checking overhead versus 100\% validity with significantly higher computational cost.

\subsubsection{2.2 Beam Search and Generation Strategies}

Beam search has been extensively studied for neural text generation, with custom scoring functions widely used in neural machine translation to incorporate length normalization and coverage penalties. These methods demonstrate that dual-objective optimization---combining model likelihood with auxiliary signals---can improve generation quality beyond pure likelihood maximization. However, prior work has not applied beam search with validity scoring to code generation, where syntax errors represent a dominant failure mode for small models. Pure beam search without validity guidance ($\alpha =1.0, \beta =0.0)$ yields error rates similar to greedy sampling (60-70\%), as explored in h-m2 ablation studies, because log-likelihood alone does not distinguish syntactically valid from invalid continuations. The contribution lies in identifying syntax validity as an effective auxiliary objective and demonstrating that small weight ($\beta =0.3)$ suffices for large error reduction (66.4\%) when combined with beam search's exploratory capacity.

\subsubsection{2.3 Type-Constrained Decoding for Code}

Prior work on constrained decoding for code has targeted type errors rather than syntax errors. Type-constrained decoding penalizes type-inconsistent generations using Mypy checking, applying soft logit penalties (typically -2.0) to tokens that would introduce type violations. Baseline experiments (h-m1 from earlier work) revealed that this approach fails: type-constrained decoding achieved 88\% syntax error rate versus 84\% baseline, because (1) the penalty weight (-2.0) is too weak to meaningfully guide generation, and (2) type errors constitute only 20\% of failures while syntax errors dominate at 70\%. Unlike type constraints which require expensive Mypy analysis at each step, syntax validation via AST parsing completes in <0.05ms, enabling real-time checking during beam search. Targeting the dominant failure mode (syntax) rather than minority errors (types) yields substantially larger improvements: 66.4\% error reduction versus h-m1's 37\% error increase.

\subsubsection{2.4 Small Model Code Generation}

Recent work on code generation has predominantly focused on large models where syntax accuracy is already high. Codex and Code Llama achieve over 90\% syntax validity at large scales (>10B parameters), making syntax errors a solved problem for these models. However, resource-constrained applications---IDE autocomplete with <100ms latency requirements, educational tools deployed on consumer hardware, internal automation with limited GPU access---necessitate smaller models ($\leq 7B$ parameters) where syntax errors remain prevalent. This work targets this regime: CodeLlama-7B exhibits 70.73\% syntax error rate on HumanEval. By reducing this to 23.78\% through lightweight validity scoring, small models become viable for applications previously restricted to expensive large models. This contribution is orthogonal to scaling-based improvements---validity scoring can be applied at any model size, though benefits diminish as baseline accuracy increases.

